\specialsection{Topological vector spaces}

When K is different from R or C, we shall call a semi-norm (resp. norm) on a vector space F over K an ultra-semi-norm (resp. an ultra-semi-norm which is a norm) on F (Esp. Vect. Top., ch. II, 2nd ed., §1). If γ is a semi-norm on F, and if x ∈ F, one will sometimes write \textbar\textbar x\textbar\textbar γ instead of γ(x).

We shall call a normable space a topological vector space over K whose topology can be defined by a single norm, and a Banach space a complete normable space \(^{1}\) . We shall call a polynormed space a topological vector space over K whose topology can be defined by a family of semi-norms; when K = R or C, this notion coincides with that of a locally convex space.

If E and F are polynormed spaces, we shall denote by \(\mathcal{L}_{m}(\mathrm{E};\mathrm{F})\) the space of continuous \(m\) -linear maps from \(\mathbf{E}^{m}\) into F, equipped with the topology of uniform convergence on the bounded subsets of \(\mathbf{E}^{m}\) ; it is a polynormed space. If moreover E is normed, and if \(\gamma\) is a continuous semi-norm on F, and if \(u\in\mathcal{L}_{m}(\mathrm{E};\mathrm{F})\) , one denotes by \(\|u\|_{\gamma}\) the lower bound of the real numbers \(a\geqslant0\) such that

\[
\left\| u (x _ {1}, \dots , x _ {m}) \right\| _ {\gamma} \leqslant a \| x _ {1} \| \dots \| x _ {m} \|
\]

for all \(x_{1}, \ldots, x_{m}\) in E.

